\textbf{Phần 1. Trắc nghiệm lựa chọn}
\setcounterpageref{ex}{0}

\begin{ex} %[3H5N32]
Mi-li-lít viết tắt là
\choice
{$\text{l}$}
{$\text{m}$}
{\True $\text{ml}$}
{$\text{mm}$}
\loigiai{
Mi-li-lít viết tắt là $\text{ml}$.\\
\textbf{Đáp án của câu hỏi là C}
}
\end{ex}

\begin{ex} %[3H5N32]
$1\text{ l}$ bằng bao nhiêu mi-li-lít?
\choice
{$10\text{ ml}$}
{\True $1000\text{ ml}$}
{$100\text{ ml}$}
{$10000\text{ ml}$}
\loigiai{
$1\text{ l} = 1000\text{ ml}$.\\
\textbf{Đáp án của câu hỏi là B}
}
\end{ex}

\begin{ex} %[3H5H32]
Tính: $350\text{ ml} + 150\text{ ml} = \dots$
\choice
{\True $500\text{ ml}$}
{$400\text{ ml}$}
{$600\text{ ml}$}
{$450\text{ ml}$}
\loigiai{
$350 + 150 = 500\text{ ml}$.\\
\textbf{Đáp án của câu hỏi là A}
}
\end{ex}

\begin{ex} %[3H5H32]
Tính: $1000\text{ ml} - 350\text{ ml} = \dots$
\choice
{$600\text{ ml}$}
{$700\text{ ml}$}
{$550\text{ ml}$}
{\True $650\text{ ml}$}
\loigiai{
$1000 - 350 = 650\text{ ml}$.\\
\textbf{Đáp án của câu hỏi là D}
}
\end{ex}

\begin{ex} %[3H5H32]
Một hộp sữa tươi có dung tích $180\text{ ml}$. Hỏi 4 hộp sữa tươi như thế có dung tích bao nhiêu mi-li-lít?
\choice
{\True $720\text{ ml}$}
{$680\text{ ml}$}
{$740\text{ ml}$}
{$800\text{ ml}$}
\loigiai{
$180 \times 4 = 720\text{ ml}$.\\
\textbf{Đáp án của câu hỏi là A}
}
\end{ex}

\begin{ex} %[3H5H32]
Có $500\text{ ml}$ nước ngọt chia đều vào 5 cốc. Mỗi cốc có số mi-li-lít nước ngọt là
\choice
{$80\text{ ml}$}
{$120\text{ ml}$}
{$150\text{ ml}$}
{\True $100\text{ ml}$}
\loigiai{
$500 : 5 = 100\text{ ml}$.\\
\textbf{Đáp án của câu hỏi là D}
}
\end{ex}

\begin{ex} %[3H5V32]
Một chai nước khoáng chứa $500\text{ ml}$, bạn Nam đã uống $220\text{ ml}$. Lượng nước còn lại trong chai là
\choice
{$260\text{ ml}$}
{\True $280\text{ ml}$}
{$300\text{ ml}$}
{$320\text{ ml}$}
\loigiai{
$500 - 220 = 280\text{ ml}$.\\
\textbf{Đáp án của câu hỏi là B}
}
\end{ex}

\begin{ex} %[3H5V32]
Một bình đựng $1\text{ l}$ nước ép. Rót ra 3 ly, mỗi ly $250\text{ ml}$. Lượng nước ép còn lại trong bình là
\choice
{$200\text{ ml}$}
{$300\text{ ml}$}
{\True $250\text{ ml}$}
{$350\text{ ml}$}
\loigiai{
Đã rót: $250 \times 3 = 750\text{ ml}$.\\
Còn lại: $1000 - 750 = 250\text{ ml}$.\\
\textbf{Đáp án của câu hỏi là C}
}
\end{ex}

\begin{ex} %[3H5V32]
Tính: $150\text{ ml} \times 6 = \dots$
\choice
{$800\text{ ml}$}
{$1000\text{ ml}$}
{$850\text{ ml}$}
{\True $900\text{ ml}$}
\loigiai{
$150 \times 6 = 900\text{ ml}$.\\
\textbf{Đáp án của câu hỏi là D}
}
\end{ex}

\begin{ex} %[3H5V32]
Đổi: $1\text{ l } 200\text{ ml} = \dots\text{ ml}$.
\choice
{$1020\text{ ml}$}
{$120\text{ ml}$}
{\True $1200\text{ ml}$}
{$1250\text{ ml}$}
\loigiai{
$1000 + 200 = 1200\text{ ml}$.\\
\textbf{Đáp án của câu hỏi là C}
}
\end{ex}

\begin{ex} %[3H5V32]
Một thìa si-rô chứa $5\text{ ml}$. Mỗi ngày em bé uống 3 thìa. Sau 4 ngày em bé uống hết số mi-li-lít si-rô là
\choice
{$50\text{ ml}$}
{\True $60\text{ ml}$}
{$70\text{ ml}$}
{$45\text{ ml}$}
\loigiai{
Mỗi ngày: $5 \times 3 = 15\text{ ml}$.\\
4 ngày: $15 \times 4 = 60\text{ ml}$.\\
\textbf{Đáp án của câu hỏi là B}
}
\end{ex}

\begin{ex} %[3H5C32]
Một can $2\text{ l}$ dầu đang đầy. Người ta rót vào 6 chai, mỗi chai $300\text{ ml}$. Lượng dầu còn lại trong can là
\choice
{\True $200\text{ ml}$}
{$150\text{ ml}$}
{$250\text{ ml}$}
{$300\text{ ml}$}
\loigiai{
$2\text{ l} = 2000\text{ ml}$.\\
Đã rót: $300 \times 6 = 1800\text{ ml}$.\\
Còn lại: $2000 - 1800 = 200\text{ ml}$.\\
\textbf{Đáp án của câu hỏi là A}
}
\end{ex}

\textbf{Phần 2. Trắc nghiệm đúng sai}
\setcounterpageref{ex}{0}

\begin{ex} %[3H5H32]
Xét tính đúng sai về đơn vị mi-li-lít:
\choiceTF
{\True $1\text{ l} = 1000\text{ ml}$.}
{\True $500\text{ ml}$ bằng nửa lít nước.}
{$1\text{ l } 50\text{ ml} = 1500\text{ ml}$.}
{\True $1\text{ l } 50\text{ ml} = 1050\text{ ml}$.}
\loigiai{
\begin{itemize}
    \item $1\text{ l} = 1000\text{ ml}$: Đúng.
    \item $500\text{ ml} = \dfrac{1}{2}\text{ l}$: Đúng.
    \item $1\text{ l } 50\text{ ml} = 1000 + 50 = 1050\text{ ml} \neq 1500\text{ ml}$. Ý 3 Sai, Ý 4 Đúng.
\end{itemize}
\begin{itemize}
    \item a) \textbf{Đúng}.
    \item b) \textbf{Đúng}.
    \item c) \textbf{Sai}.
    \item d) \textbf{Đúng}.
\end{itemize}
}
\end{ex}

\begin{ex} %[3H5V32]
Có 1 bình trà $1000\text{ ml}$. Rót đầy vào 4 cốc, mỗi cốc $200\text{ ml}$. Xét tính đúng sai:
\choiceTF
{\True Đã rót ra tất cả $800\text{ ml}$ trà.}
{\True Lượng trà còn lại trong bình là $200\text{ ml}$.}
{Lượng trà còn lại rót được thêm 2 cốc đầy nữa.}
{\True Lượng trà còn lại rót vừa đủ đúng 1 cốc đầy nữa.}
\loigiai{
\begin{itemize}
    \item Đã rót: $200 \times 4 = 800\text{ ml}$. Đúng.
    \item Còn lại: $1000 - 800 = 200\text{ ml}$. Đúng.
    \item Còn $200\text{ ml}$ chỉ rót thêm đúng 1 cốc đầy (không phải 2). Ý 3 Sai, Ý 4 Đúng.
\end{itemize}
\begin{itemize}
    \item a) \textbf{Đúng}.
    \item b) \textbf{Đúng}.
    \item c) \textbf{Sai}.
    \item d) \textbf{Đúng}.
\end{itemize}
}
\end{ex}

\textbf{Phần 3. Trả lời ngắn}
\setcounterpageref{ex}{0}

\begin{ex} %[3H5H32]
Tính: $600\text{ ml} + 400\text{ ml}$.
\shortans{1000}
\loigiai{
\textbf{Đáp số:} $1000$\\
$600 + 400 = 1000\text{ ml}$.
}
\end{ex}

\begin{ex} %[3H5H32]
Tính: $850\text{ ml} - 350\text{ ml}$.
\shortans{500}
\loigiai{
\textbf{Đáp số:} $500$\\
$850 - 350 = 500\text{ ml}$.
}
\end{ex}

\begin{ex} %[3H5V32]
Có 5 chai nước khoáng, mỗi chai chứa $200\text{ ml}$. Tính tổng dung tích theo mi-li-lít.
\shortans{1000}
\loigiai{
\textbf{Đáp số:} $1000$\\
$200 \times 5 = 1000\text{ ml}$.
}
\end{ex}

\begin{ex} %[3H5V32]
Đổi: $1\text{ l } 500\text{ ml}$ ra mi-li-lít.
\shortans{1500}
\loigiai{
\textbf{Đáp số:} $1500$\\
$1000 + 500 = 1500\text{ ml}$.
}
\end{ex}

\begin{ex} %[3H5V32]
Một bình nước có $900\text{ ml}$, rót đều vào 6 cốc. Hỏi mỗi cốc có bao nhiêu mi-li-lít nước?
\shortans{150}
\loigiai{
\textbf{Đáp số:} $150$\\
$900 : 6 = 150\text{ ml}$.
}
\end{ex}

\begin{ex} %[3H5C32]
Một chai thuốc ho $120\text{ ml}$. Mỗi lần uống $10\text{ ml}$, mỗi ngày uống 3 lần. Hỏi chai thuốc ho đó dùng được trong bao nhiêu ngày?
\shortans{4}
\loigiai{
\textbf{Đáp số:} $4$\\
Mỗi ngày uống: $10 \times 3 = 30\text{ ml}$.\\
Số ngày dùng: $120 : 30 = 4$ ngày.
}
\end{ex}
